\chapter{Całki zależne od parametru}
\label{ma1:ch:parameter-integrals}

Całka może definiować nową funkcję: jeśli $t$ jest parametrem, to wyrażenie $F(t)=\int_a^b f(x,t)\,\dd x$ zależy od $t$, mimo że zmienna całkowania $x$ znika z wyniku. Zbadamy, kiedy granicę lub pochodną względem $t$ wolno przenieść pod znak całki. Nie zakładamy znajomości teorii miary; korzystamy z całki Riemanna i całek niewłaściwych. W pierwszych czterech sekcjach przyjmujemy $a,b\in\R$, $a<b$; całki o równych lub odwróconych końcach rozumiemy zgodnie z konwencją orientacji z rozdziału \ref{ma1:ch:riemann-integral}.

\section{Całkowanie i granica na przedziale domkniętym}
\begin{twierdzenie}
\label{ma1:thm:riemann-uniform-limit}
Niech $a<b$ i niech $f_n:[a,b]\to\R$ będą całkowalne w sensie Riemanna. Jeżeli $f_n$ dążą jednostajnie do $f$, to $f$ jest całkowalna oraz
\[
 \lim_{n\to\infty}\int_a^b f_n(x)\,\dd x
 =\int_a^b f(x)\,\dd x.
\]
\end{twierdzenie}
\begin{proof}
Wybierając $n$ tak, by $|f_n-f|\le1$ na $[a,b]$, widzimy, że $f$ jest ograniczona, gdyż $f_n$ jest ograniczona. Ustalmy $\varepsilon>0$. Wybierzmy $n$ tak, by $\sup_{[a,b]}|f_n-f|<\varepsilon/(4(b-a))=:\eta$. Dla każdego podprzedziału $J$ mamy $\sup_J f\le\sup_J f_n+\eta$ i $\inf_J f\ge\inf_J f_n-\eta$. Zatem dla dowolnego podziału $P$
\[
 U(f,P)-L(f,P)\le U(f_n,P)-L(f_n,P)+2\eta(b-a).
\]
Z całkowalności $f_n$ wybieramy $P$ z różnicą sum mniejszą od $\varepsilon/2$. Wtedy różnica sum dla $f$ jest mniejsza od $\varepsilon$. Kryterium Darboux \ref{ma1:thm:darboux-criterion} daje całkowalność $f$. Z liniowości i oszacowania całki mamy
\[
 \left|\int_a^b f_n-\int_a^b f\right|
 \le (b-a)\sup_{[a,b]}|f_n-f|\longrightarrow0.
\]
\end{proof}

\begin{wniosek}
\label{ma1:cor:parameter-uniform-limit}
Niech $E\subset\R$, $t_0$ będzie punktem skupienia $E$, a $f(\cdot,t)$ będzie całkowalna na $[a,b]$ dla każdego $t\in E\setminus\{t_0\}$. Jeśli istnieje funkcja $g:[a,b]\to\R$ taka, że
\[
 \lim_{\substack{t\to t_0\\t\in E\setminus\{t_0\}}}
 \sup_{x\in[a,b]}|f(x,t)-g(x)|=0,
\]
to $g$ jest całkowalna i granicą $\int_a^b f(x,t)\,\dd x$ jest $\int_a^b g(x)\,\dd x$.
\end{wniosek}
\begin{proof}
Ponieważ $t_0$ jest punktem skupienia $E$, dla każdego $n\ge1$ wybieramy $t_n\in E$ takie, że $0<|t_n-t_0|<1/n$. Wówczas $t_n\to t_0$. Twierdzenie \ref{ma1:thm:riemann-uniform-limit} stosowane do $f(\cdot,t_n)$ daje całkowalność $g$. Następnie dla każdego $t$ stosujemy oszacowanie
\[
 \left|\int_a^b f(x,t)\,\dd x-\int_a^b g(x)\,\dd x\right|
 \le(b-a)\sup_x|f(x,t)-g(x)|.
\]
Prawa strona dąży do zera zgodnie z założeniem, co daje granicę niezależną od ciągu.
\end{proof}

\begin{przyklad}
\label{ma1:ex:parameter-moving-spike}
Sama zbieżność punktowa nie wystarcza. Dla $n\ge2$ na $[0,1]$ definiujemy
\[
 f_n(x)=
 \begin{cases}
 n^2x,&0\le x\le1/n,\\
 2n-n^2x,&1/n\le x\le2/n,\\
 0,&2/n\le x\le1.
 \end{cases}
\]
Wartości na wspólnych końcach zgadzają się, więc $f_n$ jest ciągła. Dla $x=0$ każda wartość jest zerowa, a dla $x>0$ mamy $f_n(x)=0$ od chwili, gdy $n>2/x$. Zatem $f_n\to0$ punktowo. Jednak
\[
 \int_0^1 f_n(x)\,\dd x
 =\left[\frac{n^2x^2}{2}\right]_0^{1/n}
 +\left[2nx-\frac{n^2x^2}{2}\right]_{1/n}^{2/n}
 =\frac12+\frac12=1.
\]
Maksimum wynosi $n$, więc zbieżność nie jest jednostajna. Granica całek jest różna od całki granicy.
\end{przyklad}
\begin{center}
\begin{tikzpicture}[x=6.5cm,y=.67cm,>=Stealth]
 \draw[->] (-.03,0)--(1.08,0) node[right] {$x$};
 \draw[->] (0,-.15)--(0,6.6) node[above] {$y$};
 \draw[blue!70!black,thick] (0,0)--(.5,2)--(1,0);
 \draw[teal!70!black,thick] (0,0)--(.25,4)--(.5,0);
 \draw[orange!85!black,thick] (0,0)--({1/6},6)--({1/3},0);
 \node[blue!70!black,right] at (.5,2) {$n=2$};
 \node[teal!70!black,right] at (.25,4) {$n=4$};
 \node[orange!85!black,right] at ({1/6},6) {$n=6$};
 \node[below] at (0,0) {$0$};
 \node[below] at (1,0) {$1$};
\end{tikzpicture}

\small Rysunek: trójkątne wykresy $f_n$ zwężają się, a ich wysokości rosną. Pole każdego trójkąta wynosi $1$. Dla każdego ustalonego $x$ wartość funkcji ostatecznie znika.
\end{center}

\section{Ciągłość względem parametru}
\begin{twierdzenie}
\label{ma1:thm:parameter-continuity}
Niech $I\subset\R$ będzie otwartym przedziałem. Jeżeli $f:[a,b]\times I\to\R$ jest ciągła względem pary $(x,t)$, to funkcja
\[
 F(t)=\int_a^b f(x,t)\,\dd x
\]
jest ciągła na $I$.
\end{twierdzenie}
\begin{proof}
Dla ustalonego $t$ funkcja $x\mapsto f(x,t)$ jest ciągła, więc całkowalna. Ustalmy $t_0\in I$. Wykażemy, że $f(\cdot,t)\to f(\cdot,t_0)$ jednostajnie przy $t\to t_0$. Gdyby tak nie było, istniałaby liczba $\varepsilon_0>0$ taka, że dla każdego $n\ge1$ można wybrać $t_n\in I$ z $0<|t_n-t_0|<1/n$ oraz $\sup_{x\in[a,b]}|f(x,t_n)-f(x,t_0)|\ge2\varepsilon_0$. Z definicji supremum wybieramy $x_n\in[a,b]$ z $|f(x_n,t_n)-f(x_n,t_0)|>\varepsilon_0$. Z ciągu $x_n$ wybieramy podciąg $x_{n_k}\to x_*\in[a,b]$, korzystając ze zwartości ciągowej przedziału. Wówczas obie pary $(x_{n_k},t_{n_k})$ oraz $(x_{n_k},t_0)$ dążą do $(x_*,t_0)$. Z ciągłości $f$ obie wartości dążą do $f(x_*,t_0)$, więc ich różnica dąży do zera, wbrew oszacowaniu. Otrzymaliśmy zbieżność jednostajną. Wniosek \ref{ma1:cor:parameter-uniform-limit} daje $F(t)\to F(t_0)$.
\end{proof}
Ten sam argument działa dla parametru z dowolnego podzbioru $\R$, z ciągłością względną; w punktach izolowanych ciągłość jest automatyczna. W tym kryterium zakładamy ciągłość względem obu zmiennych łącznie, a nie tylko ciągłość względem każdej z nich osobno.

\section{Różniczkowanie pod znakiem całki}
Pochodna cząstkowa $\partial_t f(x,t)$ oznacza zwykłą pochodną funkcji $s\mapsto f(x,s)$ w punkcie $t$, przy ustalonym $x$.

\begin{twierdzenie}
\label{ma1:thm:parameter-differentiation}
Niech $I$ będzie otwartym przedziałem. Jeżeli $f$ i $\partial_t f$ są ciągłe na $[a,b]\times I$, a pochodna cząstkowa istnieje w każdym punkcie, to $F(t)=\int_a^b f(x,t)\,\dd x$ jest klasy $C^1$ i
\[
 F'(t)=\int_a^b\partial_t f(x,t)\,\dd x.
\]
\end{twierdzenie}
\begin{proof}
Ustalmy $t_0\in I$. Dla $h\ne0$ takiego, że $t_0+h\in I$, mamy z liniowości
\[
 \frac{F(t_0+h)-F(t_0)}h
 =\int_a^b\frac{f(x,t_0+h)-f(x,t_0)}h\,\dd x.
\]
Ponieważ $I$ jest przedziałem, odcinek o końcach $t_0,t_0+h$ zawiera się w $I$. Funkcja ilorazu różnicowego jest ciągła względem $x$ dla każdego ustalonego $h\ne0$, więc jest całkowalna. Z twierdzenia Lagrange'a \ref{ma1:thm:lagrange}, przy ustalonym $x$, iloraz pod całką jest równy $\partial_t f(x,t_0+\theta h)$ dla pewnego $\theta\in(0,1)$ zależnego od $x,h$. Z argumentu zwartości w dowodzie twierdzenia \ref{ma1:thm:parameter-continuity}, zastosowanego do $\partial_t f$, wynika, że dla każdego $\varepsilon>0$ istnieje $\delta>0$ takie, że
\[
 |\partial_t f(x,s)-\partial_t f(x,t_0)|<\varepsilon
 \quad (x\in[a,b],\ |s-t_0|<\delta,\ s\in I).
\]
Dla $0<|h|<\delta$ oszacowanie obejmuje także $s=t_0+\theta h$, niezależnie od wyboru $\theta$. Stosując oszacowanie z dokładnością $\varepsilon/2$ zamiast $\varepsilon$, otrzymujemy supremum błędu ilorazów różnicowych nie większe niż $\varepsilon/2<\varepsilon$. Dążą one zatem jednostajnie do $\partial_t f(\cdot,t_0)$. Wniosek \ref{ma1:cor:parameter-uniform-limit} pozwala przejść do granicy pod całką. Ciągłość $F'$ wynika z twierdzenia \ref{ma1:thm:parameter-continuity}, zastosowanego do $\partial_t f$.
\end{proof}
\begin{wniosek}
\label{ma1:cor:parameter-higher-derivatives}
Niech $k$ będzie nieujemną liczbą całkowitą, $\partial_t^0f=f$, a $I$ otwartym przedziałem. Jeżeli istnieją i są ciągłe na $[a,b]\times I$ wszystkie pochodne cząstkowe $\partial_t^j f$ dla $0\le j\le k$, to $F$ jest klasy $C^k$ oraz
$F^{(j)}(t)=\int_a^b\partial_t^j f(x,t)\,\dd x$ dla $0\le j\le k$.
\end{wniosek}
\begin{proof}
Dla $j=0$ jest to definicja i twierdzenie o ciągłości. Jeżeli wzór zachodzi dla $j<k$, to stosujemy twierdzenie \ref{ma1:thm:parameter-differentiation} do $\partial_t^j f$: jej pochodną cząstkową jest $\partial_t^{j+1}f$ i obie funkcje są ciągłe. Daje to wzór i ciągłość pochodnej następnego rzędu. Indukcja kończy dowód.
\end{proof}
\begin{przyklad}
\label{ma1:ex:parameter-log}
Dla $t>-1$ definiujemy
\[
 F(t)=\int_0^1\frac{\log(1+tx)}{x}\,\dd x,
\]
przyjmując w $x=0$ wartość $t$. Dla $0\le x\le1$ i $t>-1$ mamy $1+tx\ge\min(1,1+t)>0$. Ciągłość względem pary także w $x=0$ wynika z twierdzenia Lagrange'a: dla $x>0$ i $t\ne0$ stosujemy je do logarytmu na odcinku o końcach $1,1+tx$. Iloraz jest równy $t/(1+\theta tx)$ dla pewnego $\theta\in(0,1)$; dla $t=0$ ta równość zachodzi z dowolnym $\theta\in(0,1)$. Przy $x\to0$ i $t\to t_0$ mianownik dąży do $1$, a licznik do $t_0$. Pochodna cząstkowa względem $t$ wynosi $1/(1+tx)$ także w $x=0$; jest ciągła. Zatem
\[
 F'(t)=\int_0^1\frac{\dd x}{1+tx}
 =\begin{cases}\log(1+t)/t,&t\ne0,\\1,&t=0.\end{cases}
\]
Dla $t\ne0$ korzystamy z funkcji pierwotnej $x\mapsto\log(1+tx)/t$, a dla $t=0$ całkujemy funkcję stałą $1$. Mamy również $F(0)=0$. Nie twierdzimy, że $F$ ma elementarną postać; różniczkowanie pod znakiem całki dostarcza informacji o jej pochodnej.
\end{przyklad}

\section{Reguła Leibniza dla zmiennych końców}
\begin{twierdzenie}[Reguła Leibniza]
\label{ma1:thm:parameter-leibniz}
Niech $I$ będzie otwartym przedziałem, $f,\partial_t f$ będą ciągłe na $[a,b]\times I$, a $\alpha,\beta:I\to[a,b]$ różniczkowalne. Wtedy
\[
 \begin{split}
 \frac{\dd}{\dd t}\int_{\alpha(t)}^{\beta(t)} f(x,t)\,\dd x
 ={}& f(\beta(t),t)\beta'(t)-f(\alpha(t),t)\alpha'(t)\\
 &+\int_{\alpha(t)}^{\beta(t)}\partial_t f(x,t)\,\dd x.
 \end{split}
\]
Dopuszczamy równe i odwrócone końce zgodnie z orientacją całki.
\end{twierdzenie}
\begin{proof}
Oznaczmy całkę przez $G(t)$. Ustalmy $t\in I$ i zapiszmy $\alpha_0=\alpha(t)$, $\beta_0=\beta(t)$, $\alpha_h=\alpha(t+h)$, $\beta_h=\beta(t+h)$. Addytywność daje dokładną tożsamość
\[
 \begin{split}
 G(t+h)-G(t)={}&
 \int_{\alpha_0}^{\beta_0}(f(x,t+h)-f(x,t))\,\dd x\\
 &+\int_{\beta_0}^{\beta_h}f(x,t+h)\,\dd x
 -\int_{\alpha_0}^{\alpha_h}f(x,t+h)\,\dd x.
 \end{split}
\]
Po podzieleniu przez $h$ pierwszy składnik dąży do całki po $\partial_t f$ z twierdzenia \ref{ma1:thm:parameter-differentiation} dla ustalonych końców; dla odwróconych zmieniamy znak, a dla równych oba wyrażenia są zerowe.

Od drugiego składnika podzielonego przez $h$ odejmujemy
$f(\beta_0,t)(\beta_h-\beta_0)/h$. Moduł różnicy jest nie większy niż
\[
 \left|\frac{\beta_h-\beta_0}{h}\right|
 \sup_{x\text{ między }\beta_0\text{ i }\beta_h}
       |f(x,t+h)-f(\beta_0,t)|.
\]
Pierwszy czynnik jest ograniczony przy $h\to0$, bo iloraz dąży do $\beta'(t)$. Drugi dąży do zera z ciągłości w punkcie $(\beta_0,t)$: dla każdego $x$ między końcami $|x-\beta_0|\le|\beta_h-\beta_0|\to0$. Jeśli końce są równe, różnica jest zerowa. Stąd drugi składnik po podzieleniu przez $h$ dąży do $f(\beta_0,t)\beta'(t)$. Dla trzeciego składnika podzielonego przez $h$ odejmujemy $-f(\alpha_0,t)(\alpha_h-\alpha_0)/h$. Moduł różnicy jest nie większy niż
\[
 \left|\frac{\alpha_h-\alpha_0}{h}\right|
 \sup_{x\text{ między }\alpha_0\text{ i }\alpha_h}|f(x,t+h)-f(\alpha_0,t)|.
\]
Pierwszy czynnik jest ograniczony, a drugi dąży do zera z ciągłości $f$ i $\alpha$. Stąd granicą trzeciego składnika podzielonego przez $h$ jest $-f(\alpha_0,t)\alpha'(t)$. Sumując trzy granice, otrzymujemy wzór.
\end{proof}
\begin{przyklad}
\label{ma1:ex:parameter-leibniz}
Dla $|t|<1$ rozważmy $G(t)=\int_0^t (x^2+tx)\,\dd x$ (można wybrać $[a,b]=[-1,1]$). Reguła Leibniza daje
\[
 G'(t)=2t^2+\int_0^t x\,\dd x=\frac52t^2.
\]
Bezpośrednio $G(t)=t^3/3+t^3/2=5t^3/6$, co potwierdza wzór również dla $t<0$ i w punkcie $t=0$.
\end{przyklad}

\section{Jednostajna zbieżność całek niewłaściwych}
Niech $a\in\R$, $a<b\le\infty$, $E\subset\R$ będzie niepusty, a $f(\cdot,t)$ całkowalna na każdym $[a,r]$, $a<r<b$, dla $t\in E$. Zapis $r\to b$ oznacza dojście do prawego końca od wnętrza.

\begin{definicja}
\label{ma1:def:uniform-improper}
Całka niewłaściwa $F(t)=\int_a^b f(x,t)\,\dd x$ jest zbieżna jednostajnie na $E$, jeżeli jest zbieżna dla każdego $t\in E$ i dla każdego $\varepsilon>0$ istnieje $r_0\in(a,b)$ takie, że
\[
 \left|F(t)-\int_a^r f(x,t)\,\dd x\right|<\varepsilon
 \qquad(t\in E,\ r_0<r<b).
\]
Istotne jest to, że $r_0$ nie zależy od parametru.
\end{definicja}
\begin{twierdzenie}[Jednostajne kryterium Cauchy'ego]
\label{ma1:thm:uniform-improper-cauchy}
Całka niewłaściwa jest zbieżna jednostajnie na $E$ wtedy i tylko wtedy, gdy dla każdego $\varepsilon>0$ istnieje $r_0\in(a,b)$ takie, że
\[
 \left|\int_u^v f(x,t)\,\dd x\right|<\varepsilon
 \qquad(t\in E,\ r_0<u<v<b).
\]
\end{twierdzenie}
\begin{proof}
Jeżeli całka jest zbieżna jednostajnie, wybieramy $r_0$ z definicji dla $\varepsilon/2$ i odejmujemy całki obcięte, używając nierówności trójkąta. Odwrotnie, ustalając $t$, stosujemy kryterium Cauchy'ego \ref{ma1:thm:improper-cauchy}; otrzymujemy istnienie $F(t)$. Dla zadanego $\varepsilon>0$ wybieramy wspólne $r_0$ w warunku z dokładnością $\varepsilon/2$. Dla dowolnych $t\in E$ i $r_0<u<b$ przechodzimy do granicy $v\to b$ w nierówności $|\int_u^v f|<\varepsilon/2$. Otrzymujemy
\[
 \left|F(t)-\int_a^u f(x,t)\,\dd x\right|\le\varepsilon/2<\varepsilon.
\]
To daje zbieżność jednostajną.
\end{proof}
\begin{twierdzenie}[Kryterium majoranty]
\label{ma1:thm:uniform-improper-majorant}
Jeżeli $|f(x,t)|\le m(x)$ dla $x\in[a,b)$ i $t\in E$, a $m\ge0$ jest lokalnie całkowalna i $\int_a^b m(x)\,\dd x$ jest zbieżna, to całka z $f$ jest zbieżna bezwzględnie dla każdego $t$ i jednostajnie na $E$. Także całka z $|f|$ jest zbieżna jednostajnie na $E$.
\end{twierdzenie}
\begin{proof}
Na każdym obcięciu z całkowalności modułu i porównania wynika
\[
 \left|\int_u^v f(x,t)\,\dd x\right|
 \le\int_u^v |f(x,t)|\,\dd x\le\int_u^v m(x)\,\dd x.
\]
Prawa strona jest mniejsza od $\varepsilon$ dla wszystkich $r_0<u<v<b$ dzięki kryterium Cauchy'ego dla $m$. Wybór $r_0$ jest niezależny od $t$. Stosujemy twierdzenie \ref{ma1:thm:uniform-improper-cauchy} do $f$ i $|f|$.
\end{proof}

\begin{twierdzenie}[Zamiana granicy z całką niewłaściwą]
\label{ma1:thm:improper-parameter-limit}
Niech $t_0$ będzie punktem skupienia $E$. Załóżmy, że całka z $f$ jest zbieżna jednostajnie na $E\setminus\{t_0\}$ oraz dla każdego $r\in(a,b)$ funkcje $f(\cdot,t)$ dążą jednostajnie na $[a,r]$ do tej samej funkcji $g$ przy $t\to t_0$, $t\in E\setminus\{t_0\}$. Wtedy całka niewłaściwa z $g$ jest zbieżna i
\[
 \lim_{\substack{t\to t_0\\t\in E\setminus\{t_0\}}}
 \int_a^b f(x,t)\,\dd x=\int_a^b g(x)\,\dd x.
\]
\end{twierdzenie}
\begin{proof}
Wniosek \ref{ma1:cor:parameter-uniform-limit} daje całkowalność $g$ na każdym $[a,r]$ i zbieżność całek obciętych. Ustalmy $\varepsilon>0$. Z jednostajnego kryterium Cauchy'ego wybieramy $r_0$ tak, by dla $r_0<u<v<b$ i wszystkich dopuszczalnych $t$ zachodziło $|\int_u^v f(x,t)\,\dd x|<\varepsilon/2$. Dla ustalonych $u,v$ przechodzimy do granicy $t\to t_0$ na przedziale $[u,v]$ (jednostajność wynika z jednostajności na $[a,v]$). Dostajemy $|\int_u^v g(x)\,\dd x|\le\varepsilon/2<\varepsilon$. Kryterium Cauchy'ego daje zbieżność całki z $g$.

Dla wykazania wzoru wybieramy wspólne obcięcie $r$ tak bliskie końcowi $b$, by ogon całki z $g$ oraz ogony całek z $f(\cdot,t)$ miały moduły mniejsze od $\varepsilon/3$, niezależnie od $t$. Dla ustalonego $r$ z wniosku \ref{ma1:cor:parameter-uniform-limit} wynika, że moduł różnicy całek obciętych jest mniejszy od $\varepsilon/3$ przy $t$ dostatecznie bliskim $t_0$. Nierówność trójkąta dla dwóch ogonów i różnicy całek obciętych daje moduł różnicy pełnych całek mniejszy od $\varepsilon$.
\end{proof}

\section{Ciągłość i różniczkowanie całek niewłaściwych}
\begin{twierdzenie}
\label{ma1:thm:improper-parameter-continuity}
Niech $f:[a,b)\times I\to\R$ będzie ciągła, gdzie $I$ jest otwartym przedziałem. Jeżeli całka niewłaściwa jest zbieżna jednostajnie na każdym domkniętym przedziale $K\subset I$, to $F(t)=\int_a^b f(x,t)\,\dd x$ jest ciągła na $I$.
\end{twierdzenie}
\begin{proof}
Ustalmy $t_0\in I$ i wybierzmy $K=[t_0-d,t_0+d]\subset I$ z $d>0$. Dla $\varepsilon>0$ wybieramy obcięcie $r$ tak, aby dla wszystkich $t\in K$ różnica między $F(t)$ a $F_r(t)=\int_a^r f(x,t)\,\dd x$ miała moduł mniejszy od $\varepsilon/3$. Funkcja $F_r$ jest ciągła z twierdzenia \ref{ma1:thm:parameter-continuity}. Dla $t$ dostatecznie bliskiego $t_0$ w $K$ mamy $|F_r(t)-F_r(t_0)|<\varepsilon/3$. Stąd
\[
 |F(t)-F(t_0)|\le |F(t)-F_r(t)|+|F_r(t)-F_r(t_0)|
 +|F_r(t_0)-F(t_0)|<\varepsilon.
\]
\end{proof}
W szczególności, gdy $f$ ma na każdym takim $K$ całkowalną majorantę niezależną od $t$, kryterium majoranty zapewnia hipotezę twierdzenia.

\begin{twierdzenie}
\label{ma1:thm:improper-parameter-differentiation}
Niech $f,\partial_t f$ będą ciągłe na $[a,b)\times I$, gdzie $I$ jest otwartym przedziałem. Załóżmy, że dla pewnego $t_*\in I$ całka $\int_a^b f(x,t_*)\,\dd x$ jest zbieżna oraz całka z $\partial_t f$ jest zbieżna jednostajnie na każdym domkniętym przedziale $K\subset I$. Wtedy całka z $f$ jest zbieżna dla każdego $t\in I$, jednostajnie na każdym takim $K$, a $F$ jest klasy $C^1$ i
\[
 F'(t)=\int_a^b\partial_t f(x,t)\,\dd x.
\]
\end{twierdzenie}
\begin{proof}
Dla $a<r<b$ połóżmy
$F_r(t)=\int_a^r f(x,t)\,\dd x$ oraz
$H_r(t)=\int_a^r\partial_t f(x,t)\,\dd x$.
Z twierdzenia \ref{ma1:thm:parameter-differentiation} mamy $F_r'=H_r$, a $H_r$ jest ciągła. Wzór Newtona--Leibniza względem parametru daje
\[
 F_r(t)=F_r(t_*)+\int_{t_*}^t H_r(s)\,\dd s.
\]
Oznaczmy granicę $H_r(s)$ przez $H(s)$. Z twierdzenia \ref{ma1:thm:improper-parameter-continuity} funkcja $H$ jest ciągła na $I$. Dla dowolnego domkniętego przedziału $K=[c,d]\subset I$ bierzemy $K'=[\min(c,t_*),\max(d,t_*)]\subset I$; zawieranie wynika z tego, że $I$ jest przedziałem. Wtedy
\[
 \left|\int_{t_*}^t(H_r(s)-H(s))\,\dd s\right|
 \le |t-t_*|\sup_{s\in K'}|H_r(s)-H(s)|
\]
dla $t\in K$. Niech $D=\max(|c-t_*|,|d-t_*|)$. Ponieważ $F_r(t_*)$ ma skończoną granicę $A$, otrzymujemy
\[
 \sup_{t\in K}\left|F_r(t)-\left(A+\int_{t_*}^t H(s)\,\dd s\right)\right|
 \le |F_r(t_*)-A|+D\sup_{s\in K'}|H_r(s)-H(s)|.
\]
Pierwszy składnik dąży do zera z założenia o punkcie bazowym, a drugi z jednostajnej zbieżności na $K'$. Zatem $F_r$ dąży jednostajnie na $K$ do funkcji
\[
 F(t)=A+\int_{t_*}^t H(s)\,\dd s.
\]
Dla każdego $t$ jest to granica obciętych całek, czyli szukana całka niewłaściwa. Z ciągłości $H$ i podstawowego twierdzenia rachunku całkowego \ref{ma1:thm:riemann-fundamental} otrzymujemy $F'=H$. Nie użyliśmy zamiany dwóch całek niewłaściwych.
\end{proof}

\begin{przyklad}
\label{ma1:ex:parameter-exponential}
Dla $t>0$ mamy
\[
 F(t)=\int_0^\infty e^{-tx}\,\dd x=\frac1t,
\]
ponieważ całka do $R$ wynosi $(1-e^{-tR})/t$. Na każdym $K=[c,d]\subset(0,\infty)$ z $c>0$ zachodzi
$|e^{-tx}|\le e^{-cx}$ i $|\partial_t e^{-tx}|=xe^{-tx}\le xe^{-cx}$.
Obie majoranty są całkowalne: pierwsza ma całkę $1/c$, a dla drugiej podstawienie $y=cx$ i przykład \ref{ma1:ex:improper-parts-exp} dają całkę $1/c^2$. Możemy zatem różniczkować pod znakiem całki i otrzymujemy
\[
 -\frac1{t^2}=F'(t)=-\int_0^\infty xe^{-tx}\,\dd x.
\]
Ogólniej, dla całkowitego $n\ge0$,
\[
 \int_0^\infty x^n e^{-tx}\,\dd x=\frac{n!}{t^{n+1}}.
\]
Istotnie, na obcięciu całkowanie przez części daje dla $n\ge1$
\[
 \int_0^R x^n e^{-tx}\,\dd x
 =-\frac{R^n e^{-tR}}{t}+\frac{n}{t}\int_0^R x^{n-1}e^{-tx}\,\dd x.
\]
Granica pierwszego składnika jest zerowa: ze szeregu wykładniczego $e^{tR}\ge(tR)^{n+1}/(n+1)!$. Wobec tego $0\le R^ne^{-tR}\le(n+1)!/(t^{n+1}R)$. Indukcja, poczynając od $n=0$, daje wzór. W szczególności na $K$ każda pochodna $\partial_t^n e^{-tx}=(-x)^ne^{-tx}$ ma całkowalną majorantę $x^ne^{-cx}$; wielokrotne różniczkowanie jest uzasadnione przez powtarzanie twierdzenia.
\end{przyklad}

\section{Kontrprzykłady przy braku kontroli ogonów}
\begin{przyklad}
\label{ma1:ex:improper-parameter-discontinuity}
Niech $f(x,t)=te^{-tx}$ dla $x\ge0$, $t\ge0$. Funkcja jest ciągła względem pary. Dla $t>0$ całka od $0$ do $R$ wynosi $1-e^{-tR}$, a dla $t=0$ wynosi zero. Stąd
\[
 F(t)=\begin{cases}1,&t>0,\\0,&t=0.\end{cases}
\]
Całka jest zbieżna dla każdego parametru, lecz $F$ nie jest ciągła w $0$ względem $[0,\infty)$. Dla każdego $R>0$ i $t>0$ ogon ma wartość $e^{-tR}$, a jego supremum po $0<t\le1$ wynosi $1$ (granica przy $t\to0^+$). Zbieżność nie jest jednostajna na $[0,1]$.
\end{przyklad}
\begin{przyklad}
\label{ma1:ex:improper-parameter-nondifferentiability}
Niech $f(x,t)=t^2 e^{-|t|x}$ dla $x\ge0$, $t\in\R$. Dla $t\ne0$
\[
 \partial_t f(x,t)=
 (2t-t^2x\operatorname{sgn}t)e^{-|t|x},
\]
gdzie $\operatorname{sgn}t=1$ dla $t>0$ i $-1$ dla $t<0$. W $t=0$ pochodna wynosi zero, gdyż iloraz różnicowy $h e^{-|h|x}$ dąży do zera. Pochodna jest ciągła względem pary także przy $t=0$: dla $x$ z ograniczonego przedziału jej moduł jest nie większy niż $2|t|+t^2x$. Dla $t\ne0$ całka z $f$ wynosi $t^2/|t|=|t|$, a w $0$ wynosi zero. Zatem $F(t)=|t|$ nie jest różniczkowalna w $0$.

Dla $t\ne0$ całka z $\partial_t f$ jest zbieżna bezwzględnie, ponieważ każdy składnik jest ograniczony przez stałą razy $e^{-|t|x}$ lub $xe^{-|t|x}$. Korzystając z przykładu \ref{ma1:ex:parameter-exponential} z dodatnim parametrem $|t|$, otrzymujemy
\[
 \int_0^\infty\partial_t f(x,t)\,\dd x
 =\frac{2t}{|t|}-\frac{t^2\operatorname{sgn}t}{|t|^2}
 =\operatorname{sgn}t.
\]
Dla $t=0$ ta całka wynosi zero. Jej funkcja parametru jest nieciągła w $0$. Brak jednostajnej zbieżności na każdym $[-d,d]$, $d>0$, można także wykazać bezpośrednio. Dla $t>0$ całka po $f$ do $R$ wynosi $t(1-e^{-tR})$; różniczkując ją na skończonym przedziale, dostajemy $\int_0^R\partial_t f=1-e^{-tR}+tRe^{-tR}$. Ogon całki z pochodnej wynosi zatem $(1-tR)e^{-tR}$. Dla $t=1/(2R)$ i dostatecznie dużych $R$ mamy $t\in[-d,d]$, a ogon jest równy $\tfrac12e^{-1/2}>0$, niezależnie od $R$. Sama całka z $f$ jest natomiast zbieżna jednostajnie nawet na $\R$: jej ogon dla $t\ne0$ wynosi $|t|e^{-|t|R}$, a dla $t=0$ zero. Dla $R>0$ supremum po $t$ wynosi $1/(eR)$, gdyż funkcja $s\mapsto se^{-sR}$ na $[0,\infty)$ ma pochodną $(1-sR)e^{-sR}$ i maksimum w $s=1/R$. Ciągłość obu funkcji podcałkowych i punktowa zbieżność obu całek nie wystarczają do różniczkowania.
\end{przyklad}

Na przedziałach domkniętych głównym narzędziem jest kontrola jednostajna względem zmiennej całkowania. Dla całek niewłaściwych trzeba ponadto kontrolować ogony niezależnie od parametru. Gdy całka ma kilka niewłaściwych końców, stosujemy odpowiednie kryteria osobno do każdej części, zgodnie z definicją \ref{ma1:def:improper-splitting}. Wersje oparte na teorii miary, w tym twierdzenie o zbieżności zmajoryzowanej, mają szerszy zakres i nie są potrzebne w podanych dowodach.
